%% ****** Start of file apstemplate.tex ****** %
%%
%%
%%   This file is part of the APS files in the REVTeX 4.2 distribution.
%%   Version 4.2a of REVTeX, January, 2015
%%
%%
%%   Copyright (c) 2015 The American Physical Society.
%%
%%   See the REVTeX 4 README file for restrictions and more information.
%%
%
% This is a template for producing manuscripts for use with REVTEX 4.2
% Copy this file to another name and then work on that file.
% That way, you always have this original template file to use.
%
% Group addresses by affiliation; use superscriptaddress for long
% author lists, or if there are many overlapping affiliations.
% For Phys. Rev. appearance, change preprint to twocolumn.
% Choose pra, prb, prc, prd, pre, prl, prstab, prstper, or rmp for journal
%  Add 'draft' option to mark overfull boxes with black boxes
%  Add 'showkeys' option to make keywords appear
%\documentclass[aps,prl,preprint,groupedaddress]{revtex4-2}
\documentclass[aps,prl,twocolumn,groupedaddress]{revtex4-2}
%\documentclass[aps,prl,preprint,superscriptaddress]{revtex4-2}
%\documentclass[aps,prl,reprint,groupedaddress]{revtex4-2}

% You should use BibTeX and apsrev.bst for references
% Choosing a journal automatically selects the correct APS
% BibTeX style file (bst file), so only uncomment the line
% below if necessary.
%\bibliographystyle{apsrev4-2}

\usepackage{graphicx}
\usepackage{epstopdf}
\usepackage{caption}
%\usepackage{amsmath}% http://ctan.org/pkg/amsmath
%\usepackage{amsthm}
%\usepackage{amsfonts}
%\usepackage{subfigure}
%\usepackage{hhline}
%\usepackage[miktex]{gnuplottex}
%\usepackage{xcolor}
\usepackage{amssymb}
\usepackage{amsmath}
\usepackage{color}
\usepackage{hyperref}
%\usepackage[percent]{overpic}
\usepackage{tikz}
\usepackage{mathrsfs}
\usepackage{wasysym}
\usepackage{tikz-cd}
%\usepackage{linearb}
%\usepackage{wsuipa}
\usepackage{calc}
\usepackage{tipa}

%\usepackage{stix} %\fisheye
\usepackage{stackengine,scalerel}

% so sections, subsections, etc. become numerated.
\setcounter{secnumdepth}{3}

\newcommand{\FF}{\mathbb{F}}
\newcommand{\NN}{\mathbb{N}}
\newcommand{\ZZ}{\mathbb{Z}}
\newcommand{\QQ}{\mathbb{Q}}
\newcommand{\RR}{\mathbb{R}}
\newcommand{\CC}{\mathbb{C}}
\newcommand{\II}{\mathbb{I}}
\newcommand{\GG}{\mathbb{G}}

\DeclareMathOperator*{\argmax}{arg\,max}
\DeclareMathOperator*{\argmin}{arg\,min}
\newcommand{\avrg}[1]{\left\langle #1 \right\rangle}
\newcommand{\nelta}{\bar{\delta}}
\newcommand{\bra}[1]{\left\langle #1\right|}
\newcommand{\ket}[1]{\left| #1 \right\rangle}
\newcommand{\sbra}[1]{\langle #1|}
\newcommand{\sket}[1]{| #1 \rangle}
\newcommand{\bek}[3]{\left\langle #1 \right| #2 \left| #3 \right\rangle}
\newcommand{\sbek}[3]{\langle #1 | #2 | #3 \rangle}
\newcommand{\braket}[2]{\left\langle #1 \middle| #2 \right\rangle}
\newcommand{\ketbra}[2]{\left| #1 \middle\rangle \middle\langle #2  \right|}
\newcommand{\sbraket}[2]{\langle #1 | #2 \rangle}
\newcommand{\sketbra}[2]{| #1 \rangle  \langle #2 |}
\newcommand{\norm}[1]{\left\lVert#1\right\rVert}
\newcommand{\snorm}[1]{\lVert#1\rVert}
\newcommand{\bvec}[1]{\boldsymbol{\mathsf{#1}}}
\newcommand{\bcov}[1]{\boldsymbol{#1}}
\newcommand{\bdua}[1]{\boldsymbol{\check{#1}}}
%\newcommand{\bdov}[1]{\boldsymbol{\breve{#1}}}
\newcommand{\bdov}[1]{\breve{#1}}
%\newcommand{\bten}[1]{\boldsymbol{\mathfrak{#1}}}
\newcommand{\bten}[1]{\boldsymbol{\mathfrak{#1}}}
\newcommand{\forany}{\tilde{\forall}}
\newcommand{\qed}{$\overset{\circ}{.}\;$}
\newcommand{\bigo}{\mathcal{O}}
\newcommand{\spn}{\mathrm{spn}\,}
\newcommand{\krn}{\mathrm{krn}\,}
\newcommand{\biy}{\leftrightarrow}
\newcommand{\ugh}{\mathbin{\textlinb{\BPwheel}}}

\newcommand\bigeye{\ensurestackMath{\stackinset{c}{}{c}{-.3pt}%
  {\bullet}{\scriptstyle\bigcirc}}}
\newcommand\eye{\scalerel*{\bigeye}{x}}
%\newcommand*{\fisheye}{%
%    \mathbin{%
%        \ooalign{$\circledcirc$\cr\hidewidth$\bullet$\hidewidth}%
%    }%
%}

% double angle-bracket
% https://tex.stackexchange.com/questions/79657/how-to-get-double-angle-bracket-without-using-mnsymbol-package
\makeatletter
\newsavebox{\@brx}
\newcommand{\llangle}[1][]{\savebox{\@brx}{\(\m@th{#1\langle}\)}%
  \mathopen{\copy\@brx\mkern2mu\kern-0.9\wd\@brx\usebox{\@brx}}}
\newcommand{\rrangle}[1][]{\savebox{\@brx}{\(\m@th{#1\rangle}\)}%
  \mathclose{\copy\@brx\mkern2mu\kern-0.9\wd\@brx\usebox{\@brx}}}
\makeatother

% https://tex.stackexchange.com/questions/297362/bar-through-letter-in-mathmode
%\usepackage{calc}
\newsavebox\CBox
\newcommand\hcancel[2][0.5pt]{%
  \ifmmode\sbox\CBox{$#2$}\else\sbox\CBox{#2}\fi%
  \makebox[0pt][l]{\usebox\CBox}%  
  \rule[0.77\ht\CBox-#1/2]{\wd\CBox}{#1}}
\newcommand{\cb}{\hcancel{b}}
\newcommand{\cd}{\hcancel{d}}
\newcommand{\ct}{\hcancel{\partial}}
%\newcommand{\supp}{\mathsf{supp}}
\newcommand{\supp}{\mathrm{supp}}
\newcommand{\csupp}{\overline{\mathrm{supp}}}
\newcommand{\supr}[2]{\underset{#1}{\text{supr}}}

\begin{document}

% Use the \preprint command to place your local institutional report
% number in the upper righthand corner of the title page in preprint mode.
% Multiple \preprint commands are allowed.
% Use the 'preprintnumbers' class option to override journal defaults
% to display numbers if necessary
%\preprint{}

%Title of paper
\title{
M\'etodos Matem\'aticos de la F\'isica II\\
Part I:
Structure of Banach and Hilbert spaces
}

% repeat the \author .. \affiliation  etc. as needed
% \email, \thanks, \homepage, \altaffiliation all apply to the current
% author. Explanatory text should go in the []'s, actual e-mail
% address or url should go in the {}'s for \email and \homepage.
% Please use the appropriate macro foreach each type of information

% \affiliation command applies to all authors since the last
% \affiliation command. The \affiliation command should follow the
% other information
% \affiliation can be followed by \email, \homepage, \thanks as well.
\author{Juan I. Perotti}
\email[]{juan.perotti@unc.edu.ar}
%\homepage[]{Your web page}
%\thanks{}
%\altaffiliation{}
%\affiliation{}
\affiliation{Instituto de F\'isica Enrique Gaviola (IFEG-CONICET), Ciudad Universitaria, 5000 C\'ordoba, Argentina}
\affiliation{Facultad de Matem\'atica, Astronom\'ia, F\'isica y Computaci\'on, Universidad Nacional de C\'ordoba, Ciudad Universitaria, 5000 C\'ordoba, Argentina}

% \author{Colaborador}
% \affiliation{Instituto de F\'isica Enrique Gaviola (IFEG-CONICET), Ciudad Universitaria, 5000 C\'ordoba, Argentina}
% \affiliation{Facultad de Matem\'atica, Astronom\'ia, F\'isica y Computaci\'on, Universidad Nacional de C\'ordoba, Ciudad Universitaria, 5000 C\'ordoba, Argentina}
% \email[E-mail: ]{xxx.yyy@unc.edu.ar }

%Collaboration name if desired (requires use of superscriptaddress
%option in \documentclass). \noaffiliation is required (may also be
%used with the \author command).
%\collaboration can be followed by \email, \homepage, \thanks as well.
%\collaboration{Claudio J. Tessone}
%\noaffiliation

\date{\today}

\begin{abstract}
Bla bla bla... resumen...
\end{abstract}

% insert suggested keywords - APS authors don't need to do this
%\keywords{}

%\maketitle must follow title, authors, abstract, and keywords
\maketitle
\onecolumngrid

\section{Metric space}

A set $X$ is called a {\em metric space} if it is equipped with a notion of {\em metric distance}, i.e. a function $d:X\times X\to \RR$ such that, for any three elements $x$, $y$ and $z$ in $X$, which are called {\em points}, it holds:
\begin{enumerate}
\item $d(x,y) \geq 0$, i.e. it is {\em non-negative}.
\item $d(x,y) = 0$ implies $x=y$, i.e. it is {\em positive definite}.
\item $d(x,y) = d(y,x)$, i.e. it is {\em symmetric}.
\item $d(x,z) \leq  d(x,y) + d(y,z)$, i.e. it satisfies the {\em triangle inequality}.
\end{enumerate}

{\em Example 1:} The set $\RR^n$ with the euclidean distance $d(x,y) = \sqrt{\sum_{i=1}^n (x_i-y_i)^2}$ is a metric space.

{\em Example 2:}
For any $p\geq 1$, the set 
$$
L^p(\RR^n)
=
\bigg\{
f\in C(\RR^n)
:
\int_{\RR^n}
dx\,
|f(x)|^p
<
\infty
\bigg\}
$$
is a metric space with metric distance
$$
d(f,g)
=
\bigg(
\int_{\RR^n}
dx\,
|f(x)-g(x)|^p
\bigg)^{1/p}
.
$$
At difference with previous example, this is a metric space that is associated to an infinite-dimensional vector space.

\subsection{Convergence}

We say that a sequence $\{x_j\in X:j\in \NN\}$ {\em converges} to a value $x\in X$ if, for each $\epsilon>0$, there exists an $N\in \NN$ such that $d(x_j,x)<\epsilon$ for every $j\geq N$.
It can be easily shown that, when the sequence is convergent, then the value of $x$ is unique.

The convergence of the sequence is abbreviated by $x_j\to x$ as $j\to \infty$, or even more simply by just $x_j\to x$.


\subsection{Continuous function}

The notion of convergence enables the definition of the continuity of a function.
Namely, consider two metric spaces $X$ and $Y$.
A function $f:X\to Y$ is continuous at $x\in X$ if, for every convergent sequence $x_j \to x$ it holds that $f(x_j)\to f(x)$.

\subsection{
Cauchy sequence
\label{sec:cauchy}
}

A sequence $\{x_j\in X:j\in \NN\}$ is a {\em Cauchy sequence} if, for every $\epsilon>0$, there exists an $N\in \NN$ such that $d(x_i,x_j)<\epsilon$ for every $i,j\geq N$.

{\em Excercise:} Show that a convergent sequence is a Cauchy sequence.

{\em Remark:} Not every Cauchy sequence in a metric space $X$ converges to a value in $X$.
The following example illustrates this idea.

{\em Example:}
Consider the sequence given by
$$
x_j = \sum_{i=0}^j \frac{1}{i!}.
$$
Clearly, $x_j\to e$ as $j\to \infty$, i.e. it converges to Euler's number $e=2.71828...$~.
This is a Cauchy sequence in $\QQ$ but it does not converge to a value in $\QQ$ because $e$ is irrational.
In other words, it converges to a value in $\RR$ that is not in $\QQ$.

\subsection{Complete metric space}

A metric space $X$ equipped with a metric $d$ is called {\em complete} if every Cauchy sequence in it converges to an element in it.
More formally, if $X$ is a metric space, then for any Cauchy sequence $\{x_j\in X:j\in \NN\}$, there exists an $x\in X$ such that, for every $\epsilon>0$, there exists an $N\in \NN$ such that $d(x_j,x)<\epsilon$ for every $j\geq N$.

{\em Remark:}
The example of previous section,~\ref{sec:cauchy} implies that $\QQ$ is not a complete space.

\section{Normed vector space}

A {\em normed vector space} $V$ over a field $\FF$ is a vector space equipped with a norm, which is a mapping $V\ni v\to |v|\in \RR$ that, for any $v,w\in V$ and $\alpha\in \FF$, it satisfies:
\begin{enumerate}
\item triangle inequality, $|v+w|\leq |v|+|w|$,
\item absolute homogeneity, $|\alpha v|=|\alpha||v|$, and
\item positive definiteness, $|v|=0$ implies $v=0$.
\end{enumerate}
In particular, a {\em semi-norm} is a norm without property 3.

{\em Example 1:}
In correspondence with the first example of previous section,
$$
|x| = \sqrt{\sum_{i=1}^n x_i^2}
$$
defines a norm over $\RR^n$.

{\em Example 2:}
In correspondence with the second example of previous section,
$$
|f| 
= 
\bigg(
\int_{\RR^n}
dx\,
|f(x))|^p
\bigg)^{1/p}
$$
defines a norm over $L^p(\RR^n)$.

\subsection{A normed vector space is a metric space}

A normed vector space is a metric space because a norm defines a metric distance
$$
d(v,w) = |v-w|.
$$
To see this, it is enough to check that it actually satisfies the triangle inequality, because the remaining metric properties follow directly from the norm axioms.
Namely,
$$
d(v,w) = |v-w| = |v-u+u-w| \leq |v-u|+|u-w| = d(v,u) + d(u,w)
$$
holds for any $v,u,w\in V$ due to property 1 of the norm.

{\em Remark:}
The converse is not true.
There are metric distances that do not be obtained from a norm.
For instance, given a norm $|\cdot|_a$, the expression
$$
d(x,y)
=
\frac{|x-y|_a}{1+|x-y|_a}
$$
defines a metric distance over $\RR^n$ that cannot be written as $d(x,y)=|x-y|_b$ for any norm $|\cdot|_b$.

\section{Banach space}

A complete normed vector space is called a {\em Banach space}.
In other words, as we already saw, a normed vector space is a metric space.
If additionally, such metric space is complete, meaning that all Cauchy sequences in it converge to vectors in it, then it is called a Banach space.

{\em Remark:}
Why we need Banach spaces?
Why are normed spaces not enough?
These questions will be answered in Part II but, essentially, because we want infinite series  to be convergent within the space itself.

\section{Inner product space}

Given a vector space $V$ over $\CC$, an {\em hermitian inner product} is a function $V\times V\ni (v,w) \to v\cdot w \in \CC$ satisfying, for each $v,w,u\in V$ and $\alpha,\beta\in \CC$, the following properties:
\begin{enumerate}
\item linearity of the second argument, $v\cdot(\alpha w+\beta u) = \alpha (v\cdot w) + \beta (v\cdot u)$,
\item conjugate symmetry, $v\cdot w = (w\cdot v)^*$, and
\item positive definiteness, $v\cdot v\geq 0$ and $v\cdot v=0$ only if $v=0$.
\end{enumerate}
A vector space equipped with an (hermitian) inner product is called an inner product space.

{\em Remark:} We adopt the physics convention where the inner product is linear in the second argument and antilinear in the first.

\subsection{Induced norm}

The mapping $V\ni v\to \sqrt{v\cdot v}$ is a norm called the {\em induced norm}.
Hence, an inner product space is a normed space.

\subsection{Cauchy-Schwartz inequality}

An hermitian inner product satisfies the so called Cauchy-Schwartz inequality
$$
|v\cdot w| \leq |v||w|
$$
which holds for all $v,w\in V$.

{\em Excercise:}
Prove the Cauchy-Schwartz inequality.

\section{Hilbert space}

A Hilbert space is a complete hermitian inner product space.
In other words, in a Hilbert space $H$, all Cauchy sequences converge to a vector in $H$, where the induced norm is used to determine convergence.

{\em Remark:}
Why we need Hilbert spaces?
Why are inner product spaces not enough?
As with Banach spaces, these questions will be answered in Part II These questions will be answered in Part II but, essentially, because we want linear combinations with infinite terms to be convergent within the space itself.

%\appendix
%
%\section{
%Compactness in infinite dimension
%\label{app:A}
%}
%
%% https://chatgpt.com/share/6a0c5541-7498-83e9-a94f-8ca91aa5e791
%
%The compactness of a subset $K$ of a metric space $X$ only makes sense if $X$ is finite-dimensional?
%No. Compactness is a fundamental topological notion and makes perfect sense in {\em any} metric space, including infinite-dimensional ones.
%For a subset $K \subset X$ of a metric space $(X,d)$, compactness means: 
%\begin{quote}
%Every open cover of $K$ admits a finite subcover.
%\end{quote}
%In metric spaces, this is equivalent to: 
%\begin{quote}
%Every sequence in $K$ has a convergent subsequence whose limit lies in $K$.
%\end{quote}
%This definition does not depend on dimension at all.
%What {\em does} change drastically between finite- and infinite-dimensional spaces is {\bf which sets are compact}.
%
%\subsection{Finite-dimensional case}
%
%In $\mathbb{R}^n$, the Heine–Borel theorem says:
%$$
%K \subset \mathbb{R}^n
%\text{ is compact }
%\iff
%K \text{ is closed and bounded.}
%$$
%This is a very special property of finite-dimensional normed spaces.
%
%\subsection{Infinite-dimensional case}
%
%In infinite-dimensional normed spaces, closed and bounded sets are {\em usually not compact}.
%For example, consider the Hilbert space
%$$
%\ell^2
%=
%\left\{
%x=(x_1,x_2,\dots):
%\sum_{n=1}^\infty |x_n|^2 < \infty
%\right\}
%.
%$$
%Its closed unit ball
%$$
%B=\{x\in \ell^2:|x|\le 1\}
%$$
%is closed and bounded, but {\bf not compact}.
%A standard proof uses the canonical basis vectors
%$$
%e_1=(1,0,0,\dots),\quad
%e_2=(0,1,0,\dots),\quad \dots
%$$
%which satisfy
%$$
%|e_n-e_m|=\sqrt{2}
%\quad (n\neq m).
%$$
%Hence no subsequence can be Cauchy, so no subsequence converges. Therefore $B$ is not compact.
%
%\subsection{What replaces Heine–Borel?}
%
%In metric spaces, compactness is closely related to {\bf total boundedness}.
%
%For metric spaces:
%
%$$
%K \text{ compact}
%\iff
%K \text{ complete and totally bounded}.
%$$
%
%Total boundedness is stronger than ordinary boundedness.
%A set is totally bounded if:
%\begin{quote}
%For every $\varepsilon>0$, it can be covered by finitely many balls of radius $\varepsilon$.
%\end{quote}
%In finite dimensions, boundedness already implies total boundedness. In infinite dimensions, it generally does not.
%
%\subsection{Intuition}
%
%Infinite-dimensional balls contain ``too many independent directions.''
%The sequence $e_n$ in $\ell^2$ keeps jumping into new orthogonal directions and never accumulates anywhere.
%That phenomenon cannot happen inside bounded subsets of finite-dimensional spaces.
%So compactness absolutely makes sense in infinite-dimensional spaces — it simply becomes much more restrictive.


\end{document}
%
% ****** End of file apstemplate.tex ******


